% =============================================================
% บทที่ 2 - ทฤษฎีและงานวิจัยที่เกี่ยวข้อง
% =============================================================

% -----------------------------------------------------------
\section{ทฤษฎีที่เกี่ยวข้อง}
% -----------------------------------------------------------

\subsection{เว็บสแครปปิง (Web Scraping) และการดึงข้อมูลอัตโนมัติ}

เว็บสแครปปิงคือกระบวนการดึงข้อมูลจากเว็บไซต์โดยอัตโนมัติ แทนการคัดลอกข้อมูลด้วยมือ โดยทั่วไปจะส่งคำขอ HTTP ไปยังเว็บเซิร์ฟเวอร์แล้ววิเคราะห์โครงสร้าง HTML ของหน้าเว็บเพื่อสกัดข้อมูลที่ต้องการ~\cite{barbaresi-2021-trafilatura} เครื่องมือที่นิยมใช้ ได้แก่ ไลบรารี BeautifulSoup สำหรับวิเคราะห์และค้นหาองค์ประกอบในเอกสาร HTML ด้วย CSS Selector และไลบรารี httpx สำหรับส่งคำขอ HTTP แบบอะซิงโครนัส อย่างไรก็ตาม เว็บไซต์บางแห่งแสดงผลเนื้อหาผ่าน JavaScript ต้องใช้เครื่องมือจำลองเบราว์เซอร์ เช่น Playwright เพื่อให้ได้ข้อมูลที่แสดงผลสมบูรณ์

นอกจากนี้ เว็บไซต์หลายแห่งยังเผยแพร่ข่าวผ่าน RSS Feed ซึ่งเป็นรูปแบบ XML มาตรฐานที่ใช้แจกจ่ายเนื้อหาที่อัปเดตเป็นประจำ ทำให้สามารถดึงหัวข้อข่าว ลิงก์ และบทสรุปสั้น ๆ ได้โดยไม่ต้องวิเคราะห์โครงสร้าง HTML โดยตรง~\cite{rss2009} ในระบบนี้ใช้ทั้งสองเทคนิคร่วมกันตามความเหมาะสมของแต่ละสำนักข่าว

\subsection{การสรุปข้อความด้วยโมเดลภาษาขนาดใหญ่}

การสรุปข้อความอัตโนมัติ (Text Summarization) ในทางวิทยาการคอมพิวเตอร์สามารถจำแนกออกได้เป็น 2 แนวทางหลัก ได้แก่ การสรุปแบบดึงประโยค (Extractive Summarization) ซึ่งอาศัยอัลกอริทึมทางสถิติหรือโครงสร้างกราฟ เช่น TextRank ในการคัดเลือกประโยคที่มีอยู่เดิมในบทความมาต่อเรียงกัน ซึ่งแม้จะมีความเร็วสูงแต่มีข้อจำกัดด้านความต่อเนื่องของบริบทและความเยิ่นเย้อของภาษา และการสรุปแบบสร้างข้อความใหม่ (Abstractive Summarization) ซึ่งอาศัยการทำความเข้าใจความหมายและบริบทของเนื้อหาทั้งหมด แล้วจึงสังเคราะห์เป็นข้อความสรุปชุดใหม่ด้วยภาษาที่เป็นธรรมชาติและสละสลวย~\cite{allahyari2017}

สำหรับการประมวลผลสรุปเนื้อหาข่าวในโครงงานนี้ ได้เลือกประยุกต์ใช้โมเดลภาษาขนาดใหญ่ (Large Language Model: LLM) ตระกูล Gemini ผ่านบริการ KKU Gen.AI API \cite{gemini2024} เพื่อดำเนินการสรุปแบบสร้างข้อความใหม่ โดยพิจารณาจากขีดความสามารถในการทำความเข้าใจภาษาไทยที่ซับซ้อน สามารถสกัดใจความสำคัญของบทความข่าวขนาดยาวและเรียบเรียงให้อยู่ในรูปแบบที่กระชับและถูกต้องตามหลักไวยากรณ์ นอกจากนี้ ตัวแบบยังรองรับการกำหนดรูปแบบผลลัพธ์แบบมีโครงสร้าง (Structured JSON Output) ตาม Schema ที่กำหนดไว้อย่างเคร่งครัด ทำให้ระบบสามารถสกัดบทสรุปความยาว 2--3 ประโยค ประเด็นสำคัญ การวิเคราะห์อารมณ์ของข่าวสาร (Sentiment Analysis) และคำสำคัญ (Keywords) ได้พร้อมกันในการเรียกใช้เพียงครั้งเดียว อีกทั้งยังรองรับการประมวลผลเนื้อหาหลายภาษา (Multilingual Support) โดยสรุปผลตามภาษาต้นฉบับของบทความ และการเชื่อมต่อประมวลผลผ่านบริการ API ภายนอกยังช่วยลดข้อจำกัดด้านทรัพยากรฮาร์ดแวร์ ทำให้เซิร์ฟเวอร์ของระบบไม่จำเป็นต้องพึ่งพาหน่วยประมวลผลกราฟิก (GPU) สมรรถนะสูง

\subsection{การจำแนกข้อความภาษาไทยด้วย LinearSVC และ TF-IDF}

การจำแนกข้อความ (Text Classification) มีขั้นตอนสำคัญในการแปลงข้อความตัวอักษรให้อยู่ในรูปเวกเตอร์ตัวเลข (Feature Extraction) และการใช้โมเดลการเรียนรู้ของเครื่องในการแบ่งแยกหมวดหมู่

\subsubsection{ทฤษฎี TF-IDF (Term Frequency-Inverse Document Frequency)}
TF-IDF เป็นค่าน้ำหนักเชิงสถิติที่ใช้วัดความสำคัญของคำ $t$ ในเอกสาร $d$ เทียบกับคลังเอกสารทั้งหมด $D$:
\begin{equation}
\text{TF-IDF}(t, d, D) = \text{TF}(t, d) \times \text{IDF}(t, D)
\end{equation}
โดยที่ $\text{TF}(t, d)$ คือความถี่ของคำ $t$ ในเอกสาร $d$ และ $\text{IDF}(t, D) = \log\left(\frac{1 + |D|}{1 + |\{d \in D : t \in d\}|}\right) + 1$ ช่วยลดน้ำหนักของคำที่ปรากฏบ่อยทั่วไปในทุกข่าว และเพิ่มน้ำหนักให้คำเฉพาะกลุ่มหมวดหมู่

\subsubsection{ทฤษฎี Support Vector Machine (LinearSVC)}
Linear Support Vector Classifier (LinearSVC) เป็นโมเดลที่มุ่งหาเส้นแบ่งระนาบหลายมิติ (Hyperplane) ที่มีระยะห่าง (Margin) จากจุดข้อมูลของแต่ละคลาสมากที่สุด:
\begin{equation}
\min_{w, b, \xi} \frac{1}{2} \|w\|^2 + C \sum_{i=1}^n \xi_i \quad \text{subject to } y_i (w^T x_i + b) \ge 1 - \xi_i, \quad \xi_i \ge 0
\end{equation}

การนำโมเดล LinearSVC มาใช้ร่วมกับคุณลักษณะ TF-IDF มีความเหมาะสมอย่างยิ่งต่อการจำแนกข้อความ เนื่องจากเวกเตอร์ที่ได้จากกระบวนการ TF-IDF มักมีมิติของข้อมูลสูงและมีลักษณะเบาบาง (High-dimensional Sparse Data) ซึ่ง Linear Kernel ของ SVM มีคุณสมบัติเด่นในการแบ่งกลุ่มข้อมูลลักษณะนี้ได้อย่างแม่นยำ พร้อมทั้งมีกลไก $L_2$ Regularization ที่ช่วยควบคุมการเกิด Overfitting ได้ดีกว่าโมเดล Decision Tree หรือ Naive Bayes นอกจากนี้ยังใช้เวลาในการทำนายผลรวดเร็วระดับต่ำกว่า 5 มิลลิวินาทีต่อบทความ ทำให้ระบบสามารถจำแนกหมวดหมู่ข่าวปริมาณมากในเบื้องหลังได้โดยไม่ส่งผลกระทบต่อความเร็วการให้บริการ และตัวโมเดลยังมีขนาดกะทัดรัด ประหยัดหน่วยความจำของระบบได้อย่างมีประสิทธิภาพ

\subsection{การจำแนกเชิงลึกด้วยโมเดล Transformer (WangchanBERTa)}

WangchanBERTa \cite{wangchanberta2021} เป็นโมเดลภาษาขนาดใหญ่ที่ใช้สถาปัตยกรรม RoBERTa (Robustly Optimized BERT Approach) ผ่านการฝึกฝนล่วงหน้า (Pre-trained) บนคลังข้อมูลภาษาไทยขนาดใหญ่กว่า 78 GB โดยโครงงานได้นำโมเดลดังกล่าวมาใช้เป็นส่วนเสริมสำหรับการประมวลผลข้อความเชิงลึก เนื่องจากกลไก Multi-head Self-Attention ช่วยให้โมเดลสามารถเข้าใจบริบทสองทิศทางได้อย่างแม่นยำ จึงสามารถจัดการกับคำพ้องรูป คำกำกวม หรือคำสแลงในพาดหัวข่าวที่โมเดลกลุ่ม Bag-of-Words ทั่วไปไม่สามารถตรวจจับได้ นอกจากนี้ ยังใช้เป็นเกณฑ์เปรียบเทียบมาตรฐาน (Baseline Benchmark) ในการประเมินความคุ้มค่าระหว่างความแม่นยำเชิงสัมพัทธ์กับระยะเวลาการประมวลผล

\subsection{การจัดการแคชหลายระดับและระบบคุกกี้เพื่อความเป็นส่วนตัว}

\subsubsection{ระบบแคชหลายระดับ (Multi-tier Cache Architecture)}
เพื่อตอบสนองต่อผู้ใช้งานได้อย่างรวดเร็วและลดภาระการเชื่อมต่อภายนอก ระบบได้ออกแบบสถาปัตยกรรมแคชสองระดับ ได้แก่ แคชระดับหน่วยความจำ (In-Memory LRU Cache) ผ่านคลาส \texttt{AsyncInMemoryCache} สำหรับจัดเก็บผลลัพธ์การสกัดข่าวและผลสรุปจากโมเดลภาษาขนาดใหญ่ไว้ในหน่วยความจำ RAM ทำให้คำขอเรียกดูข่าวเดิมสามารถตอบสนองได้ในเวลาต่ำกว่า 5 มิลลิวินาที และช่วยป้องกันการเรียกใช้ API ซ้ำซ้อน และแคชระดับโพรโทคอล (HTTP 304 Conditional Cache) ซึ่งอาศัยกลไก ETag และส่วนหัว \texttt{If-Modified-Since} ในการตรวจสอบสถานะการอัปเดตของหน้าเว็บต้นทาง หากไม่มีการเปลี่ยนแปลงข้อมูล เซิร์ฟเวอร์จะตอบกลับด้วยสถานะ HTTP 304 Not Modified ทำให้ลดปริมาณการส่งผ่านข้อมูลในเครือข่ายและป้องกันการถูกระงับการเข้าถึงจากเว็บไซต์ต้นทาง

\subsubsection{ระบบคุกกี้และความเป็นส่วนตัว (Privacy-Preserving Cookies \& PDPA)}
ระบบใช้คุกกี้ในการจัดเก็บค่ากำหนดเฉพาะบุคคล (Personalize Cookies) และบันทึกคะแนนปฏิสัมพันธ์ของผู้ใช้ (Algorithm Cookies) สำหรับจัดอันดับความนิยมของข่าว (Trending Algorithm) โดยคำนึงถึงพระราชบัญญัติคุ้มครองข้อมูลส่วนบุคคล (PDPA) ด้วยการไม่จัดเก็บข้อมูลระบุตัวตนส่วนบุคคล (Non-PII Aggregate Metrics) และมีแถบแจ้งเตือนขอความยินยอมการใช้คุกกี้ (Cookie Consent Banner)

% -----------------------------------------------------------
\section{งานวิจัยที่เกี่ยวข้อง}
% -----------------------------------------------------------

ในส่วนนี้จะนำเสนองานวิจัยหรือระบบที่มีลักษณะคล้ายคลึงกับโครงงานนี้ เพื่อเป็นแนวทางในการพัฒนา โดยพิจารณาครอบคลุมในด้านการสกัดเนื้อหาจากเว็บ การสรุปข้อความ และการจำแนกหมวดหมู่ข่าว

\subsection{งานสกัดเนื้อหาบทความจากเว็บ}

Barbaresi \cite{barbaresi-2021-trafilatura} ได้พัฒนาไลบรารี Trafilatura สำหรับสกัดข้อความหลักออกจากหน้าเว็บโดยอัตโนมัติ โดยไม่ขึ้นกับเว็บไซต์ใดเว็บไซต์หนึ่ง ไลบรารีนี้รองรับการจัดรูปแบบ HTML ที่หลากหลายและคืนข้อความในรูปแบบ Markdown ซึ่งเป็นประโยชน์โดยตรงต่อการนำเนื้อหาไปใช้กับโมเดลภาษาขนาดใหญ่ อย่างไรก็ตาม วิธีนี้มีข้อจำกัดเมื่อเนื้อหาถูกเรนเดอร์ด้วย JavaScript บางกรณี ทำให้ระบบนี้ต้องเสริมด้วยการจำลองเบราว์เซอร์ด้วย Playwright สำหรับเว็บไซต์กลุ่มดังกล่าว

\subsection{งานสรุปข้อความอัตโนมัติ}

Allahyari และคณะ \cite{allahyari2017} ได้สำรวจเทคนิคการสรุปข้อความอย่างเป็นระบบทั้งแบบ Extract และ Abstractive และชี้ให้เห็นว่าการสรุปแบบ Abstractive ที่ใช้การเรียนรู้เชิงลึกให้ผลลัพธ์ที่เป็นธรรมชาติกว่าอย่างมีนัยสำคัญ แต่ต้องใช้ทรัพยากรในการคำนวณสูง ด้วยความก้าวหน้าของโมเดลภาษาขนาดใหญ่ \cite{gemini2024} การสรุปแบบ Abstractive จึงเข้าถึงได้ง่ายผ่าน API โดยไม่ต้องฝึกโมเดลเอง โครงงานนี้จึงเลือกใช้แนวทางดังกล่าวเพื่อสรุปข่าวเป็นภาษาเดียวกับต้นฉบับโดยไม่ต้องใช้ทรัพยากรฝึกโมเดลของตนเอง

\subsection{งานจำแนกหมวดหมู่ข่าวภาษาไทย}

งานวิจัยด้านการจำแนกข้อความภาษาไทยมักประสบปัญหาการเว้นวรรคคำ จึงนิยมใช้เครื่องมือตัดคำภาษาไทย เช่น PyThaiNLP \cite{pythainlp2021} ร่วมกับโมเดลแมชชีนเลิร์นนิง เช่น SVM กับคุณลักษณะ TF-IDF ซึ่งเป็นแนวทางที่มีประสิทธิภาพและใช้ทรัพยากรต่ำ โครงงานนี้ประยุกต์ใช้แนวทางไฮบริด โดยใช้กฎคำสำคัญที่มีน้ำหนักและลำดับความสำคัญของ URL สำหรับการจำแนกแบบรวดเร็ว และใช้โมเดล SVM เป็นตัวตรวจสอบบริบทสำหรับหมวดหมู่ที่กฎคำสำคัญมีความแม่นยำไม่เพียงพอ

\subsection{งานวิจัยและทฤษฎีการจัดอันดับกระแสข่าวสาร (Trending News Ranking Algorithms)}

การจัดอันดับข่าวสารที่ได้รับความนิยม (Trending News) จำเป็นต้องผสานปัจจัยทั้งด้านพฤติกรรมผู้ใช้ ความสดใหม่ของเวลา และความน่าเชื่อถือของเนื้อหาเข้าด้วยกัน โดยได้รับแรงบันดาลใจจากอัลกอริทึมจัดอันดับของ Reddit \cite{salihefendic2010, swartz2005} และ Hacker News \cite{graham2009} ซึ่งประยุกต์ใช้ฟังก์ชันลอการิทึม $(\ln)$ ในการแปลงคะแนนปฏิสัมพันธ์ เพื่อลดผลกระทบของการปั่นยอดเข้าชม (Anti-Click Spam) และสะท้อนหลักการประโยชน์ส่วนเพิ่มที่ลดน้อยลง (Law of Diminishing Marginal Returns) ร่วมกับแบบจำลองการจัดอันดับข่าวที่คำนึงถึงมิติเวลาของ Dong และคณะ \cite{dong2010} ซึ่งใช้ฟังก์ชันการลดทอนแบบครึ่งชีวิต (Half-Life Exponential Decay) เพื่อให้คะแนนความนิยมของข่าวเก่าลดลงตามระยะเวลาที่ผ่านไป เปิดโอกาสให้ข่าวใหม่ขึ้นมาแสดงแทนที่ นอกจากนี้ Das และคณะ \cite{das2007} ในระบบ Google News ยังแสดงให้เห็นว่าการรวมกลุ่มข่าวประเด็นเดียวกันที่รายงานโดยสำนักข่าวอิสระหลายแห่ง (Story Clustering) เป็นตัวบ่งชี้ที่มีความแม่นยำสูงว่าประเด็นดังกล่าวเป็นเหตุการณ์สำคัญระดับสาธารณะ (Major Breaking Event)

\subsection{งานวิจัยด้านการจัดกลุ่มข่าวหลายสำนักและการหาฉันทามติของหมวดหมู่ (Multi-Source Story Clustering \& Category Consensus)}

ในระบบรวบรวมข่าวสารจากหลากหลายแหล่ง การรวมข่าวที่รายงานเหตุการณ์เดียวกันเข้าเป็นกลุ่มประเด็นเดียว (Story Cluster) เป็นกระบวนการสำคัญตามแนวคิด Topic Detection and Tracking (TDT) โดยในงานวิจัยด้านการประมวลผลข่าวสารพบว่า การเปรียบเทียบความคล้ายคลึงจากหัวข้อข่าวเพียงอย่างเดียวมักก่อให้เกิดความคลาดเคลื่อนได้ง่าย ทั้งในลักษณะพลาดการรวมข่าวเรื่องเดียวกัน (False Negatives) อันเนื่องมาจากสำนักข่าวต่างสำนักมีมุมมองการพาดหัวข่าวที่แตกต่างกัน (Journalistic Framing Variations) หรือในลักษณะรวมข่าวคนละเหตุการณ์เข้าด้วยกัน (False Positives) จากการที่ข่าวประจำวันมีรูปแบบพาดหัวที่ใช้คำศัพท์คล้ายกันแม้เป็นคนละสถานที่

เพื่อแก้ไขข้อจำกัดดังกล่าว งานวิจัยสมัยใหม่จึงนิยมผสานคุณลักษณะหลายมิติร่วมกัน ได้แก่ ข้อความพาดหัว (Title) สาระสำคัญย่อ (Summary) เอนทิตีชื่อเฉพาะบุคคลและสถานที่ (Named Entities) และคำสำคัญเฉพาะทาง (Domain Keywords) ร่วมกับการวัดความคล้ายคลึงเชิงความหมาย (Semantic Similarity) เพื่อยืนยันว่าเป็นเหตุการณ์เดียวกันอย่างแท้จริง นอกจากนี้ สำหรับปัญหาความไม่สอดคล้องของหมวดหมู่ข่าวอันเกิดจากโครงสร้าง URL และ Tag ที่แตกต่างกันในแต่ละสำนักข่าว การประยุกต์ใช้ทฤษฎีการลงมติเสียงข้างมาก (Majority Voting) หรือการทำนายซ้ำจากเนื้อหารวมของกลุ่ม จึงเป็นแนวทางมาตรฐานในการระบุหมวดหมู่ที่เป็นเอกภาพและแม่นยำที่สุดสำหรับกลุ่มข่าวดังกล่าว

\subsection{สรุปผลการเปรียบเทียบงานวิจัยที่เกี่ยวข้อง}

จากการทบทวนวรรณกรรมข้างต้น สามารถสรุปผลการเปรียบเทียบคุณสมบัติได้ดังตารางที่ \ref{tab:compare_research}

\begin{table}[H]
\centering
\caption{เปรียบเทียบงานวิจัยที่เกี่ยวข้อง}
\label{tab:compare_research}
\resizebox{\textwidth}{!}{%
\begin{tabular}{|l|c|c|c|c|c|c|}
\hline
\textbf{งานวิจัย} & \textbf{สกัดเนื้อหา} & \textbf{สรุป LLM} & \textbf{จำแนกหมวด} & \textbf{Real-time} & \textbf{แคช/คุกกี้} & \textbf{ระบบแนะนำ} \\
\hline
Trafilatura \cite{barbaresi-2021-trafilatura} & \cmark & -- & -- & -- & -- & -- \\ \hline
งานสรุปข้อความ \cite{allahyari2017} & -- & \cmark & -- & -- & -- & -- \\ \hline
งานจำแนกข้อความภาษาไทย \cite{pythainlp2021} & -- & -- & \cmark & -- & -- & -- \\ \hline
\textbf{โครงงานนี้} & \cmark & \cmark & \cmark & \cmark & \cmark & \cmark \\
\hline
\end{tabular}%
}
\end{table}

จากตารางที่ \ref{tab:compare_research} พบว่ายังไม่มีงานใดที่ผสมผสานการสกัดเนื้อหา การสรุปด้วยโมเดลภาษาขนาดใหญ่ การจำแนกหมวดหมู่ การแสดงผลแบบเรียลไทม์ การจัดการแคชและคุกกี้ และอัลกอริทึมแนะนำข่าวในระบบเดียวครบทุกด้าน โครงงานนี้จึงมุ่งพัฒนาเว็บแอปพลิเคชันที่รวบรวมคุณสมบัติทั้งหมดไว้ในระบบเดียวอย่างครบถ้วน

% -----------------------------------------------------------
\section{เหตุผลความจำเป็นในการเลือกเทคโนโลยีและเครื่องมือ (Tech Stack Rationales)}
% -----------------------------------------------------------

การเลือกชุดเทคโนโลยี (Tech Stack) ในโครงงานนี้คำนึงถึงประสิทธิภาพการทำงาน ความเข้ากันได้ของระบบ และมาตรฐานทางวิศวกรรมซอฟต์แวร์ โดยมีเหตุผลในการเลือกใช้แต่ละเทคโนโลยี ดังนี้

\subsection{FastAPI}
FastAPI \cite{fastapi2024} ถูกเลือกเป็นเฟรมเวิร์คหลักของส่วน Backend เนื่องจากทำงานบนสถาปัตยกรรมแบบ Asynchronous (ASGI บน Starlette) ซึ่งเหมาะอย่างยิ่งสำหรับงานที่เป็น Non-blocking I/O เช่น การรัน Background Scraper พร้อมกันหลายแหล่งข่าว การเชื่อมต่อ WebSocket แบบเรียลไทม์ และการเรียก API ภายนอก นอกจากนี้ยังมีความเร็วในการประมวลผลเทียบเท่าภาษา Go หรือ Node.js และสร้างเอกสาร OpenAPI Specification อัตโนมัติ

\subsection{Pydantic}
Pydantic \cite{pydantic2024} ถูกนำมาใช้ในการควบคุมและตรวจสอบชนิดข้อมูล (Data Validation \& Serialization) ที่ระดับ Runtime โดยใช้ Type Hints ของภาษา Python การใช้ Pydantic ช่วยรับประกันว่าข้อมูลข่าวที่สกัดได้จากเว็บและข้อมูลผลลัพธ์จาก LLM จะถูกต้องตรงตาม Schema ที่กำหนดไว้เสมอ ป้องกันปัญหาข้อมูลสูญหายหรือผิดรูปแบบก่อนจัดเก็บลงฐานข้อมูล

\subsection{Playwright (Chromium Headless Browser)}
Playwright \cite{playwright2024} ถูกเลือกใช้สำหรับการดึงข้อมูลหน้าเว็บที่เรนเดอร์ด้วย JavaScript (Single Page Application เช่น React/Vue) และเว็บไซต์ที่มีกลไกป้องกันบอต เช่น Cloudflare Challenge โดย Playwright สามารถจำลองเบราว์เซอร์จริงได้อย่างสมบูรณ์ จัดการ Cookies และ Browser Context แยกกันอย่างอิสระ มีความเสถียรและทำงานได้เร็วกว่า Selenium

\subsection{PyThaiNLP}
PyThaiNLP \cite{pythainlp2021} เป็นเครื่องมือประมวลผลภาษาธรรมชาติภาษาไทยมาตรฐาน โดยระบบเลือกใช้อัลกอริทึมตัดคำ \texttt{newmm} (Dictionary-based Maximal Matching ผสานกับ Thai Character Clusters: TCC) ร่วมกับชุดคำหยุดภาษาไทย (Thai Stop words) เพื่อเตรียมข้อความก่อนนำไปสร้างเวกเตอร์ TF-IDF

\subsection{Socket.IO}
Socket.IO \cite{socketio2024} ถูกเลือกสำหรับการสื่อสารแบบ Real-time ระหว่างเซิร์ฟเวอร์กับเว็บเบราว์เซอร์ เนื่องจากมีฟีเจอร์การเชื่อมต่อใหม่โดยอัตโนมัติ (Auto-reconnection), ระบบ Fallback เป็น HTTP Long-polling กรณีที่เครือข่ายไม่อนุญาตให้ใช้ WebSocket และระบบจัดการ Event-driven architecture ที่เรียบง่ายและเสถียร

\subsection{ระบบจำลองเบราว์เซอร์แยกส่วนและเทคโนโลยี Obscura (Isolated Obscura CDP Engine)}
Obscura (\texttt{h4ckf0r0day/obscura}) ถูกนำมาใช้เป็นเอนจินจำลองเบราว์เซอร์แบบแยกส่วน (Isolated Headless Browser Container) โดยสื่อสารผ่านโปรโตคอล Chrome DevTools Protocol (CDP) ผ่านการเชื่อมต่อ WebSocket (\texttt{ws://obscura:9222}) ซึ่งมีประโยชน์เชิงวิศวกรรมที่สำคัญในการแยกภาระทรัพยากร (Resource Isolation) โดยการประมวลผลเรนเดอร์หน้าเว็บข่าวที่ซับซ้อนจะถูกแยกไปทำงานในคอนเทนเนอร์ของ Obscura โดยไม่รบกวนหน่วยความจำและ CPU ของ FastAPI Backend ช่วยป้องกันปัญหาหน่วยความจำไม่เพียงพอภายใต้ทรัพยากรของเครื่องแม่ข่าย AWS EC2 \texttt{c7i-flex.large} นอกจากนี้ Obscura ยังได้รับการปรับแต่งโครงสร้าง Browser Fingerprinting และ Canvas/WebGL เพื่อลดโอกาสการถูกบล็อกคำขอจากเว็บไซต์ข่าวที่มีระบบตรวจจับบอตขั้นสูง ตลอดจนคลาส \texttt{BrowserService} ยังได้รับการออกแบบให้มีกลไกสำรองการทำงานแบบอัตโนมัติ (Graceful Fallback) ซึ่งหากคอนเทนเนอร์ภายนอกไม่พร้อมใช้งาน ระบบจะสลับไปรัน Local Playwright Chromium Engine ภายในตัวโดยอัตโนมัติ ทำให้ระบบมีความต่อเนื่องในการให้บริการ ทั้งนี้ ในระบบปัจจุบัน Obscura ถูกนำมาใช้เพื่อรองรับงานดึงและสกัดข้อมูลข่าวสารเบื้องหลัง (Background Web Scraping \& Extraction) ภายในระบบ Container เท่านั้น โดยยังไม่ได้นำไปเชื่อมต่อใช้งานกับระบบ In-App Browser หรือส่วนแสดงผลหน้าเว็บฝั่งไคลเอนต์

\subsection{ระบบจัดเก็บข้อมูลตามหลัก SOLID/GRASP (NewsRepositoryPort)}
การออกแบบระบบจัดเก็บข้อมูลใช้ Port and Adapter Pattern ผ่านโปรโตคอล \texttt{NewsRepositoryPort} ทำให้โค้ดส่วนบริการไม่ผูกติดกับเทคโนโลยีฐานข้อมูลใดฐานข้อมูลหนึ่ง โดยระบบเริ่มต้นพัฒนาด้วย \texttt{FileNewsRepository} ที่จัดเก็บเป็นไฟล์ JSON เพื่อความสะดวกรวดเร็วในการติดตั้งและทดสอบ และสามารถสลับไปใช้ฐานข้อมูล SQL หรือ NoSQL ได้ทันทีในอนาคตโดยไม่ต้องแก้ไข Business Logic

\subsection{ระบบแคชหลายระดับ (Multi-tier Cache Strategy)}
การเลือกใช้กลยุทธ์แคชสองระดับระหว่าง \texttt{AsyncInMemoryCache} (LRU Cache ใน RAM) และ HTTP 304 Conditional Cache ช่วยตัดปัญหาการดึงข้อมูลและประมวลผลสรุปเนื้อหาซ้ำซ้อนอย่างสิ้นเชิง ส่งผลให้สามารถลดค่าใช้จ่ายในการเรียกใช้ LLM API ภายนอกได้มากกว่า 70\% และเพิ่มความเร็วในการตอบสนองคำขอเรียกดูข่าวของผู้ใช้ได้อย่างมีประสิทธิภาพสูงสุด
